% Part: first-order-logic
% Chapter: proof-systems
% Section: tableaux

\documentclass[../../../include/open-logic-section]{subfiles}

\begin{document}

\iftag{FOL}
      {\olfileid{fol}{prf}{tab}}
      {\olfileid{pl}{prf}{tab}}

\olsection{\usetoken{P}{tableau}}

در حالی که بسیاری از دستگاه‌های !!{derivation} با آرایش‌هایی از !!{sentence}s
کار می‌کنند، !!{tableau}s با !!{signed formula}s کار می‌کنند. هر
!!^a{signed formula} زوجی است متشکل از یک علامت ارزش صدق
($\True$ یا $\False$) و !!^a{sentence}:
\[
\sFmla{\True}{!A} \text{ یا } \sFmla{\False}{!A}.
\]
هر !!^a{tableau} از !!{signed formula}s تشکیل می‌شود که در درختی با انشعاب
رو به پایین آرایش یافته‌اند. تابلو با شماری \emph{فرض} آغاز می‌شود و با
!!{signed formula}s ادامه می‌یابد که از اعمال یکی از قواعد استنتاج بر یکی
از !!{signed formula}s بالای آن حاصل شده‌اند. هر قاعده به ما امکان می‌دهد یک
یا چند !!{signed formula} را به انتهای یک شاخه بیفزاییم، یا دو
!!{signed formula} را در کنار هم بیفزاییم---در حالت اخیر، شاخه به دو شاخه
تقسیم می‌شود و دو !!{signed
  formula} افزوده‌شده انتهای آن دو شاخه را می‌سازند.

قاعده‌ای که بر یک !!{signed formula} مرکب اعمال شود، به این می‌انجامد که
!!{signed formula}s افزوده شوند؛ این‌ها زیر!!{formula}s بی‌واسطهٔ آن‌اند. قواعد
به صورت جفت می‌آیند، یک قاعده برای هر یک از دو علامت. برای نمونه، قاعدهٔ
$\TRule{\True}{\land}$ بر $\sFmla{\True}{!A \land !B}$ اعمال می‌شود و
افزودن هر دو !!{signed formula} یعنی $\sFmla{\True}{!A}$ و
~$\sFmla{\True}{!B}$ را به انتهای هر شاخه‌ای که
$\sFmla{\True}{!A \land !B}$ را دربر دارد مجاز می‌سازد؛ و قاعدهٔ
$\TRule{\False}{\land}$ اجازه می‌دهد شاخه با افزودن
$\sFmla{\False}{!A}$ و $\sFmla{\False}{!B}$ در کنار هم منشعب شود. هر
!!^a{tableau} بسته است هرگاه تک‌تک شاخه‌هایش یک جفت متناظر از
!!{signed formula}s، یعنی $\sFmla{\True}{!A}$ و
$\sFmla{\False}{!A}$، را دربر داشته باشند.

رابطهٔ $\Proves$ مبتنی بر !!{tableau}s چنین تعریف می‌شود:
$\Gamma \Proves !A$ اگر و تنها اگر مجموعهٔ متناهی‌ای چون
~$\Gamma_0 = \{!B_1, \dots, !B_n\} \subseteq \Gamma$ وجود داشته باشد
که برای فرض‌های
\[
\{\sFmla{\False}{!A}, \sFmla{\True}{!B_1}, \dots, \sFmla{\True}{!B_n}\}
\]
!!^a{tableau}ی بسته‌ای وجود داشته باشد. برای نمونه، !!^a{tableau}ی بستهٔ زیر
نشان می‌دهد که $\Proves (!A \land !B) \lif !A$:
\begin{oltableau}
  [\sFmla{\False}{(\formula{A} \land \formula{B}) \lif \formula{A}}, just = \TAss
    [\sFmla{\True}{\formula{A} \land \formula{B}}, just = {\TRule{\False}{\lif}[1]}
      [\sFmla{\False}{\formula{A}}, just = {\TRule{\False}{\lif}[1]}
        [\sFmla{\True}{\formula{A}}, just = {\TRule{\True}{\land}[2]}
          [\sFmla{\True}{\formula{B}}, just = {\TRule{\True}{\land}[2]}, close]
        ]
      ]
    ]
  ]
\end{oltableau}

مجموعهٔ $\Gamma$ در حساب !!{tableau} ناسازگار است هرگاه برای فرض‌های
\[
\{\sFmla{\True}{!B_1}, \dots, \sFmla{\True}{!B_n}\}
\]
برای برخی $!B_i \in \Gamma$، !!^a{tableau}ی بسته‌ای وجود داشته باشد.

!!^{tableau}s را اورت بث و یاکو هینتیکا در دههٔ ۱۹۵۰، مستقل از یکدیگر،
ابداع کردند و ریموند اسمولیان آن‌ها را ساده و رایج کرد. کاربردشان بسیار
آسان است، زیرا ساختن !!^a{tableau} روندی بسیار نظام‌مند است. به سبب ماهیت
نظام‌مند !!{tableau}s، پیاده‌سازی رایانه‌ای آن‌ها نیز آسان است. با این حال،
خواندن !!^a{tableau} غالباً دشوار است و پیوند آن با برهان‌ها همیشه به‌آسانی
دیده نمی‌شود. این رویکرد همچنین کاملاً عام است و بسیاری از منطق‌های گوناگون
دستگاه‌های !!{tableau} دارند. !!^{tableau}s همچنین به ما کمک می‌کنند
!!{structure}sیی بیابیم که (مجموعه‌هایی از) !!{sentence}s داده‌شده را ارضا
می‌کنند: اگر مجموعه ارضاپذیر باشد، !!^a{tableau}ی بسته نخواهد داشت؛ یعنی هر
!!^a{tableau} شاخه‌ای باز خواهد داشت. !!{structure} ارضاکننده را می‌توان از
یک شاخهٔ باز «خواند»، مشروط بر اینکه هر قاعده‌ای که اعمالش ممکن است بر آن
شاخه اعمال شده باشد. پیوند بسیار نزدیکی نیز با حساب دنباله‌ای وجود دارد:
در اصل، هر !!^a{tableau}ی بسته !!^a{derivation}ی فشرده‌ای در حساب دنباله‌ای
است که وارونه نوشته شده است.

\end{document}
